\chapter{Facturation par école}
\label{ch:facturation}

Depuis la v2, la \textbf{préparation de facturation} est \textbf{contextualisée par école}. Vous travaillez mois par mois sur les sessions d'un seul établissement à la fois.

\section{Accès}

\begin{workflowstep}[Ouvrir la facturation]
  \texttt{/home} → cliquez sur une école → descendez à la section \textbf{Facturation} ou utilisez l'ancre \texttt{\#billing} dans l'URL.
\end{workflowstep}

\screenshot[keepaspectratio]{09-school-billing.png}{Section facturation sur la fiche école — navigation mensuelle}

\section{Éléments de l'écran}

\begin{table}[htbp]
  \centering
  \caption{Blocs de la section facturation}
  \begin{tabular}{@{}lp{8.5cm}@{}}
    \toprule
    \textbf{Bloc} & \textbf{Rôle} \\
    \midrule
    Sélecteur de période & Mois précédent/suivant, liste des mois \\
    Bascule \menu{Par date} & Regroupement chronologique ou par cours \\
    Tableau des sessions & Groupe, horaire, heures, montant HT, montant TTC, n° facture \\
    Totaux & Heures et montants HT/TTC par cours et par école \\
    Actions (Édition) & Rattacher une facture existante, créer une facture \\
    \bottomrule
  \end{tabular}
\end{table}

\section{Navigation mensuelle}

Les variables de session \texttt{billing\_year} et \texttt{current\_month} pilotent l'affichage. L'année de facturation est distincte de l'année académique (\texttt{current\_year}).

Chaque ligne d'intervention affiche le montant \textbf{HT} (durée réelle × tarif × multiplicateur facturable) et le \textbf{TTC} (HT ×\,1{,}2, arrondi à 2 décimales). Les totaux cours / école rappellent HT et TTC.

\begin{workflowstep}[Changer de mois]
  Utilisez les flèches ou la liste déroulante des mois. Les totaux et sessions se mettent à jour instantanément.
\end{workflowstep}

\section{Saut aux sessions non facturées}

Le bouton \menu{Sauter aux sessions non facturées} recherche en arrière le dernier mois comportant au moins une session sans facture.

\begin{itemize}[leftmargin=*]
  \item Utile pour rattraper une facturation en retard.
  \item \textbf{Désactivé} s'il n'existe aucun mois antérieur non facturé.
\end{itemize}

\section{Créer ou rattacher une facture}

En mode \textbf{Édition}~:

\begin{enumerate}[label=\textbf{\arabic*.}]
  \item Sélectionnez les sessions du mois à inclure (selon la vue par cours ou par date).
  \item \textbf{Rattacher} une facture déjà émise via \menu{Assigner une facture}.
  \item Ou \textbf{créer} une nouvelle facture~: numéro, date, montants HT/TTC, PDF généré.
\end{enumerate}

Les factures créées apparaissent ensuite dans \menu{Trésorerie → Factures}.

\section{Indicateurs sur la liste des écoles}

Sur \texttt{/home}, chaque ligne affiche~:

\begin{itemize}[leftmargin=*]
  \item \textbf{Facturé TTC}~: somme des factures de l'année courante.
  \item \textbf{Non facturé TTC}~: sessions sans \texttt{invoice\_id}, montant majoré de TVA (×\,1{,}2).
\end{itemize}

Ces chiffres utilisent les mêmes formules que la section facturation (\texttt{Tools::planningGain}, durée réelle de session).

\section{Anciennes URLs}

Les routes \texttt{/billing*} redirigent vers \texttt{school.show\#billing} si une école est en session, sinon vers la liste des écoles.

\begin{astuce}
  Marquez une facture comme payée via le \textbf{rapprochement bancaire} (chapitre~\ref{ch:tresorerie}) plutôt que manuellement, pour conserver la cohérence date de paiement / opération bancaire.
\end{astuce}
